We describe each method's allocation for 2020 Census data,
summarising the properties developed in J.1--J.4.
Table~\ref{tab:2020seats} shows the seat allocation for all 50 states
under all five methods; the three states where any method disagrees
with Huntington-Hill are highlighted.

\subsection{Huntington-Hill (Equal Proportions)}

Huntington-Hill assigns the $s$-th seat to state $i$ if
$p_i / \sqrt{(s-1)s}$ is the highest priority value among all states.
Starting from 1 seat per state (50 seats assigned), seats 51 through 435
are awarded sequentially by priority.
The 2020 apportionment under this method matches the Census Bureau's
official result: California 52, Texas 38, Florida 28, New York 26,
\ldots, Wyoming 1.
The last seat (seat 435) went to New York; Montana narrowly missed
a second seat (J.1 documents the priority margin).

\textit{Optimality property}: Huntington-Hill minimises, over all pairs
of states $(i, j)$, the maximum relative difference
$(p_i/s_i) / (p_j/s_j)$ in per-capita representation
\citep{huntington1928method}.

\subsection{Webster (Major Fractions)}

Webster adjusts a common divisor $d$ until $\sum_i [p_i/d]_{0.5} = 435$,
where $[\cdot]_{0.5}$ denotes rounding to the nearest integer
(rounding at 0.5).
Under 2020 data, Webster produces a seat allocation identical to
Huntington-Hill for 49 states; Minnesota retains its 8th seat under
Webster while losing it to another state under Huntington-Hill
(see J.2 for the exact computation).

\textit{Optimality property}: Webster minimises the sum of squared
deviations $\sum_i (s_i - q_i)^2$ where $q_i = p_i \times 435 / N$
is the exact quota \citep{balinski1982fair}.

\subsection{Adams (Smallest Divisors)}

Adams assigns each state $\lceil p_i/d \rceil$ seats for a divisor $d$
chosen such that the total is 435.
Because ceiling always rounds up, every state receives at least
$\lceil q_i \rceil$ seats --- meaning no state is below its exact
proportional entitlement.
Under 2020 data, Adams gives Rhode Island, Hawaii, and Montana one extra
seat each relative to Huntington-Hill; large states correspondingly lose
seats.

\textit{Optimality property}: Adams maximises the minimum per-capita
representation $\min_i (p_i/s_i)$ across all states
\citep{balinski1982fair}.

\subsection{Jefferson (Greatest Divisors / D'Hondt)}

Jefferson assigns each state $\lfloor p_i/d \rfloor$ seats, subject to
a floor of 1.
Floor rounding systematically underrepresents small states:
Wyoming's exact quota of $\approx 0.758$ is rounded to 1 only due to the
constitutional floor, not by Jefferson's rounding rule.
Jefferson is the dominant method in European and Latin American
proportional representation systems, where it is called d'Hondt
\citep{lijphart1986degrees}.
Under 2020 data, Jefferson would give Texas 39 and California 53 seats
while several small states would be reduced to the constitutional
minimum of 1.

\textit{Equivalence}: Jefferson and d'Hondt are mathematically identical;
the priority value $P_i(s) = p_i/s$ is the same formula
regardless of whether $p_i$ represents population or votes (J.4).

\subsection{Hamilton (Largest Remainders)}

Hamilton computes exact quotas $q_i = p_i \times 435 / N$,
assigns lower quotas $\lfloor q_i \rfloor$, and distributes the
remaining $435 - \sum_i \lfloor q_i \rfloor$ seats to states
with the largest fractional parts.
Hamilton satisfies the quota rule (every state is within 1 of its
exact quota) but is susceptible to all three classical paradoxes.
Under 2020 data, Hamilton produces the same result as Huntington-Hill
for all 50 states --- the methods agree for this particular census ---
but they can disagree in other years and especially for smaller House sizes.

\begin{table}[t]
\centering
\small
\caption{2020 Census apportionment: selected states under all five methods.
  States marked $*$ differ from Huntington-Hill under at least one method.
  Full 50-state table appears in the Appendix.}
\label{tab:2020seats}
\begin{tabular}{lrrrrrr}
\toprule
State & Pop.\ (2020) & HH & Webster & Adams & Jefferson & Hamilton \\
\midrule
California    & 39,538,223 & 52 & 52 & 51 & 53 & 52 \\
Texas         & 29,145,505 & 38 & 38 & 37 & 39 & 38 \\
Florida       & 21,538,187 & 28 & 28 & 27 & 28 & 28 \\
New York      & 20,201,249 & 26 & 26 & 26 & 26 & 26 \\
North Carolina &10,439,388 & 14 & 14 & 14 & 14 & 14 \\
Minnesota$^*$ &  5,706,494 &  8 &  8 &  8 &  7 &  8 \\
Rhode Island$^*$& 1,097,379 & 2 &  2 &  3 &  1 &  2 \\
Montana$^*$   &  1,084,225 &  2 &  2 &  3 &  1 &  2 \\
Hawaii$^*$    &  1,455,271 &  2 &  2 &  3 &  2 &  2 \\
Wyoming       &    576,851 &  1 &  1 &  1 &  1 &  1 \\
Vermont       &    643,077 &  1 &  1 &  1 &  1 &  1 \\
Alaska        &    733,391 &  1 &  1 &  2 &  1 &  1 \\
\bottomrule
\end{tabular}
\begin{tablenotes}
\small
\item HH = Huntington-Hill (current U.S.\ law).
  All methods give 435 total seats and a minimum of 1 per state.
  The Adams column reflects the small-state-favoring ceiling rounding;
  the Jefferson column reflects large-state-favoring floor rounding.
\end{tablenotes}
\end{table}

\subsection{Where Methods Diverge}

For the 2020 Census, the five methods produce identical allocations for
at least 45 of the 50 states.
Divergence concentrates near the rounding threshold for small states
and near integer-quota boundaries for mid-size states.
Three patterns emerge:
\begin{enumerate}
  \item \textbf{Small-state sensitivity}: Adams gives Alaska, Montana,
    Rhode Island, and Hawaii an extra seat relative to HH;
    Jefferson would reduce Minnesota and potentially Rhode Island
    relative to HH.
  \item \textbf{Large-state offsetting}: Any extra seats given to small
    states under Adams must be taken from large states (CA, TX, FL)
    because the total is fixed at 435.
  \item \textbf{Near-threshold states}: Minnesota at the 2020 census
    sits very close to the boundary between 7 and 8 seats;
    it receives 8 under HH, Webster, Adams, and Hamilton,
    but would receive 7 under Jefferson.
\end{enumerate}
These findings underscore that the choice of apportionment method is
not merely academic: it determines the political representation of
specific states for a decade.
